\textbf{Nombre de régimes.} Sur toute la période d'apprentissage, le BIC de l'EM est minimal à $K = 2$ (BIC = 4494,1, tableau~\ref{tab:prix_reference}) ; ce $K$ sert de $K_{\text{réf}}$ et de modèle générateur du jeu semi-synthétique.

\textbf{Ajustements sur toute la période d'apprentissage} (tableau~\ref{tab:prix_complet}). EM-K : $\hat K$ = 2, log-vraisemblance de test -4,290 ; EM-BIC : $\hat K$ = 2, log-vraisemblance de test -4,290 ; EM-ICL : $\hat K$ = 2, log-vraisemblance de test -4,290 ; MAP-EM : $\hat K$ = 7, log-vraisemblance de test -4,266 ; MAP-cal : $\hat K$ = 4, log-vraisemblance de test -4,255 ; VB-emp : $\hat K$ = 4, log-vraisemblance de test -4,318 (prédictive -4,308) ; VB-cal : $\hat K$ = 8, log-vraisemblance de test -4,238 (prédictive -4,249).

\textbf{Régime de tension et alertes} (tableau~\ref{tab:prix_alerte}). EM-BIC, apprentissage : régime de tension de poids 0,105, d'écarts moyens 0,112 (maïs), 0,008 (riz), -0,023 (niébé) ; 7,8\,\% des observations en alerte ; largeur moyenne de l'intervalle à 90 \% de la probabilité d'alerte 0,946 ; EM-BIC, test : régime de tension de poids 0,105, d'écarts moyens 0,112 (maïs), 0,008 (riz), -0,023 (niébé) ; 7,7\,\% des observations en alerte ; largeur moyenne de l'intervalle à 90 \% de la probabilité d'alerte 0,921 ; VB-emp, apprentissage : régime de tension de poids 0,145, d'écarts moyens 0,045 (maïs), 0,048 (riz), 0,150 (niébé) ; 11,2\,\% des observations en alerte ; largeur moyenne de l'intervalle à 90 \% de la probabilité d'alerte 0,075 ; VB-emp, test : régime de tension de poids 0,145, d'écarts moyens 0,045 (maïs), 0,048 (riz), 0,150 (niébé) ; 12,2\,\% des observations en alerte ; largeur moyenne de l'intervalle à 90 \% de la probabilité d'alerte 0,093 ; les deux méthodes classent de la même façon (probabilité supérieure ou non à 0,5) 83,8\,\% des observations.
